\section{Exhibit A-1: calculus as a stable closure}
\label{sec:exhibit-calculus}

This exhibit instantiates the \SBT closure lens in the most familiar mathematical setting:
how the differential and integral calculus appear as the stable closure of discrete protocols under refinement.
Our emphasis is not on definitions but on \emph{feasibility}:
which compressions remain coherent as the stage parameter refines.

\subsection{Discrete substrate and staging}
Fix a step size $h>0$ and consider functions on a one-dimensional grid (or, more generally, any set with a notion of translation by $h$).
The discrete substrate supplies:
\begin{itemize}
  \item a \emph{shift} operator $T_h$ defined by $(T_h f)(x)=f(x+h)$;
  \item a \emph{difference} operator $\Delta_h := T_h-I$;
  \item the \emph{scaled difference} (difference quotient)
  \[
    \delta_h f(x)\;:=\;\frac{f(x+h)-f(x)}{h},
  \]
  which is the simplest instance of \Pone\ operator rewriting: the $h^{-1}$ rescaling that keeps the update comparable across stages.
\end{itemize}

In \SBT terms, $h\downarrow 0$ is \Pfour\ staging.
Any proposed macro-law (``derivative'', ``generator'', ``flow'') is a packaged object only if its induced defect ledger remains feasible under refinement (\Psix\ accounting), and only if it respects the coherence constraints imposed at the packaged layer (\Ptwo\ constraints).

\subsection{The accounting ledger: exact discrete Leibniz identity}
At the discrete level, the difference quotient is not a derivation:
multiplicative coherence fails by a remainder term that can be written \emph{exactly}.
This remainder is the canonical \Psix\ ledger term for the passage from differences to derivatives.

\begin{lemma}[Finite-difference Leibniz identity]\label{lem:finite-diff-leibniz}
Let $F$ be a field, let $h\in F$ with $h\neq 0$, and let $f,g:F\to F$.
Define $\delta_h f(x):=(f(x+h)-f(x))/h$.
Then for all $x\in F$,
\[
\delta_h(fg)(x)
\;=\;
\delta_h f(x)\, g(x)
\;+
\;f(x)\, \delta_h g(x)
\;+
\;h\,\delta_h f(x)\,\delta_h g(x).
\]
\end{lemma}

\begin{remark}[Machine-checked anchor]\label{rem:lean-leibniz}
Lemma~\ref{lem:finite-diff-leibniz} is formalized in Lean/mathlib in this repository as
\texttt{Diff.delta\_mul\_leibniz} (file \texttt{lean/Diff/FiniteDifference.lean}) \cite{moura2021lean4,mathlib2020lean}.
\end{remark}

The key point for closure is not the identity itself but the structure of the ledger term:
\[
\underbrace{h\,\delta_h f\,\delta_h g}_{\text{ledger term}}
\]
is explicitly proportional to the stage parameter $h$.
Thus, in any refinement regime where $\delta_h f$ and $\delta_h g$ remain bounded (a feasibility condition measured by \Psix),
the ledger term vanishes as $h\downarrow 0$ and multiplicative coherence emerges at the packaged layer.

\subsection{Packaging: quotient by vanishing defect yields derivations}
The previous subsection suggests a general closure template.
A discrete operator family $L_h$ (or a protocol that composes many micro-steps) is admissible as a ``law'' in the packaged layer only after specifying:
\begin{enumerate}
  \item a staging parameter $h$ (\Pfour);
  \item a ledger that tracks closure defects, such as Leibniz residuals, consistency across refinement, and stability bounds (\Psix);
  \item an identification rule ``equal up to vanishing defect as $h\downarrow 0$'' that packages families into an equivalence class (\Pfive).
\end{enumerate}

Demanding multiplicative coherence is a \Ptwo\ constraint.
For difference quotients it is satisfied only \emph{approximately} at finite $h$ and becomes exact in the packaged limit because the ledger term in Lemma~\ref{lem:finite-diff-leibniz} vanishes.
The packaged limit of feasible difference-quotient families therefore satisfies the Leibniz rule:
\[
D(fg) \;=\; (Df)\,g + f\,(Dg),
\]
i.e.\ it is a \emph{derivation}.
In one dimension, further invariance/normalization choices single out the usual derivative, but the closure point is more general:
\emph{derivations are what survives after enforcing multiplicative coherence under refinement.}

\subsection{Algebraic uniqueness anchor on polynomials}
The previous discussion is analytic and asymptotic.
To anchor the ``derivation as closure fixed point'' motif in a fully algebraic setting, we use polynomial rings where derivations have a clean universal property.

\begin{theorem}[Derivations on $R\lbrack X\rbrack$ are determined by $D(X)$]\label{thm:derivation-determined-by-X}
Let $R$ be a commutative semiring and let $A:=R[X]$ be the polynomial ring.
An $R$-derivation $D:\mathrm{Derivation}(R,A,A)$ is uniquely determined by its value on the indeterminate $X\in A$.
Equivalently, evaluation on $X$ induces an $R$-linear equivalence
\[
\mathrm{Derivation}(R,R[X],R[X]) \;\simeq_{\ell[R]}\; R[X],
\qquad
D \longmapsto D(X),
\]
with inverse given (schematically) by ``multiply the formal derivative by a polynomial''.
\end{theorem}

\begin{remark}[Machine-checked anchors]\label{rem:lean-derivation-poly}
Theorem~\ref{thm:derivation-determined-by-X} is formalized in Lean/mathlib in this repository as:
\path{Derivations.derivation_ext_X} (uniqueness) and
\path{Derivations.polynomialDerivationEquiv} (explicit $R$-linear equivalence),
in file \path{lean/Derivations/Polynomial.lean}.
The Lean type of the equivalence is
\path{Derivation R (Polynomial R) (Polynomial R)} $\simeq_{\ell[R]}$ \path{Polynomial R}.
The evaluation-at-$X$ interface is exposed by
\path{Derivations.polynomialDerivationEquiv_apply_X} and
\path{Derivations.polynomialDerivationEquiv_symm_apply_X}.
These are re-exported/instantiated from mathlib with small glue lemmas. \cite{moura2021lean4,mathlib2020lean}
\end{remark}

Classically, this computation is expressed via Kähler differentials, e.g.\ $\Omega_{R[X]/R} \cong R[X]\cdot dX$; see \cite{StacksProject00RM}.

For the purposes of this paper, the takeaway is structural:
once the Leibniz rule is imposed as a coherence constraint, the space of admissible operators collapses to a small, rigid family.
In polynomial land this rigidity is literal and exact (determined by $D(X)$);
in analytic settings it is mediated by the defect ledger and the chosen packaging rule.

\subsection{Falsification-first diagnostics: stencil selection under feasibility}
To complement the algebraic anchors, we implemented a numerical experiment that samples a population of candidate operators and filters them by feasibility gates.
We sample local linear stencil operators and apply gates that instantiate \Pfour--\Psix:
stability under refinement, decreasing route mismatch, and (optionally) decreasing multiplicative (Leibniz) defect.

\subsection{Numerical protocol details (stencils)}
We emphasize that the stencil experiments are \emph{diagnostics}, not a uniqueness theorem.
A candidate local operator is a clipped-boundary stencil of half-width $m$,
\[
  (L f)_i := \sum_{j=-m}^{m} c_j\, f_{i+j},
  \qquad i=m,\dots,N-m-1,
\]
evaluated only on interior points to avoid boundary-condition artifacts.
To compare across refinements $N,2N,4N$ (step sizes $h, h/2, h/4$), we align outputs by restricting the finer-grid interior outputs to even-index centers.
The refinement-consistency defects $E_{01},E_{12}$ used in the feasibility gate are relative $\ell^2$ mismatches on this aligned interior, maximized over a fixed smooth test family (polynomials plus $\sin$ and $\exp$).

The ``order'' label $k\in\{0,1,2\}$ is chosen by best-fit against analytic templates on the finest grid:
after the renormalization $h^{-k}$ (an instance of \Pone\ operator rewriting), we select the $k$ whose output is closest (up to an overall scalar) to $f$, $f'$, or $f''$ on the test family.

In the recorded runs, the baseline grid size is $N_0=\StencilNZero$ and the Leibniz-gated experiment uses half-width $m=\StencilMLeibniz$ with gate thresholds $D_2<\StencilGateDTwo$ and ratio gate $\mathrm{gate\_ratio}=\StencilGateRatio$.
For standard background on finite-difference stencils, consistency, and stability, see \cite{LeVeque2007}.

\paragraph{Disclosure of the \Pone\ step in the Leibniz-gated run.}
In the Leibniz-gated experiment we additionally apply a moment-normalization step to candidate coefficient vectors $(c_j)$ before auditing them.
Concretely, we project to a ``first-order'' class by enforcing the discrete moment constraints
$\sum_j c_j=0$ and $\sum_j j\,c_j = 1$, and we optionally suppress low-order contamination by setting higher moments (e.g.\ $\sum_j j^2 c_j$, $\sum_j j^3 c_j$) near zero.
In \SBT terms this is not the conclusion but the \emph{operator rewrite} (\Pone) that makes families comparable across scale and avoids trivial $k=0$ closures.
The empirical point of Table~\ref{tab:stencil-leibniz} is therefore conditional:
\emph{within the normalized first-order class, imposing a shrinking Leibniz defect selects derivative-like closures.}

Two regimes illustrate the role of constraints.

\paragraph{Stability-only filtering.}
If we impose only refinement stability/consistency without demanding multiplicative coherence, then the surviving operators overwhelmingly behave as order-$0$ closures (approximately scalar multiplication).
In the current run, \StencilSurvivorsTotal\ out of \StencilGeneratedTotal\ sampled stencils survive under the stability gate, and the survivors are classified as best-fit order $k=0$
(Table~\ref{tab:stencil-stability}).

\paragraph{Leibniz-gated filtering.}
If we additionally demand that the Leibniz defect improves with refinement (a direct numerical analog of Lemma~\ref{lem:finite-diff-leibniz} having a vanishing ledger term), then the surviving operators collapse to derivative-like (order-$1$) closures.
In the current run, the Leibniz-gated filter produces \StencilLeibnizSurvivorsTotal\ survivors and all survivors are classified as best-fit order $k=1$
(Table~\ref{tab:stencil-leibniz}).
Figure~\ref{fig:stencil-leibniz} visualizes the selection outcome for the recorded run.

These experiments do not prove uniqueness of the derivative in full generality.
They serve as \emph{diagnostics} aligned with the \SBT program:
when a closure is feasible, defects shrink under refinement; when it is not, defects persist or blow up.
Here, enforcing multiplicative coherence is the decisive feasibility pressure that selects derivation-like operators from a large population of local discrete rules.

% --- Auto-generated artifacts (tables/figure) ---
% Auto-generated by scripts/export_results_tex.py from notes/*_last_run.json
\begin{table}[h]
\centering
\caption{Stencil flow stability-only summary}
\label{tab:stencil-stability}
\begin{tabular}{l|l|l}
\hline
 generated & survivors & best-fit counts \\
\hline
 1000 & 913 & k=0: 913; k=1: 0; k=2: 0 \\
\hline
\end{tabular}
\end{table}

% Auto-generated by scripts/export_results_tex.py from notes/*_last_run.json
\begin{table}[h]
\centering
\caption{Leibniz-gated stencil selection}
\label{tab:stencil-leibniz}
\begin{tabular}{l|l}
\hline
 survivors & best-fit counts \\
\hline
 50 & k=0: 0; k=1: 50; k=2: 0 \\
\hline
\end{tabular}
\end{table}

% Auto-generated by scripts/export_results_tex.py from notes/*_last_run.json
\begin{table}[h]
\centering
\caption{Leibniz-gated top survivors (diagnostics)}
\label{tab:stencil-leibniz-top}
\begin{tabular}{l|l|l|l}
\hline
 rank & D2 & F1 & \ensuremath{E_{12,1}} \\
\hline
 1 & \ensuremath{1.68144\times 10^{-6}} & \ensuremath{7.19072\times 10^{-7}} & \ensuremath{5.92487\times 10^{-6}} \\
\hline
 2 & \ensuremath{1.8767\times 10^{-6}} & \ensuremath{8.04778\times 10^{-7}} & \ensuremath{6.6106\times 10^{-6}} \\
\hline
 3 & \ensuremath{2.05543\times 10^{-6}} & \ensuremath{8.80601\times 10^{-7}} & \ensuremath{7.24186\times 10^{-6}} \\
\hline
 4 & \ensuremath{2.19982\times 10^{-6}} & \ensuremath{9.44249\times 10^{-7}} & \ensuremath{7.74517\times 10^{-6}} \\
\hline
 5 & \ensuremath{2.20896\times 10^{-6}} & \ensuremath{9.45855\times 10^{-7}} & \ensuremath{7.77949\times 10^{-6}} \\
\hline
 6 & \ensuremath{2.24716\times 10^{-6}} & \ensuremath{7.73372\times 10^{-7}} & \ensuremath{6.41209\times 10^{-6}} \\
\hline
 7 & \ensuremath{2.92045\times 10^{-6}} & \ensuremath{1.66745\times 10^{-7}} & \ensuremath{1.404\times 10^{-6}} \\
\hline
 8 & \ensuremath{3.00861\times 10^{-6}} & \ensuremath{1.29407\times 10^{-6}} & \ensuremath{1.05915\times 10^{-5}} \\
\hline
 9 & \ensuremath{3.11535\times 10^{-6}} & \ensuremath{8.37018\times 10^{-7}} & \ensuremath{6.81121\times 10^{-6}} \\
\hline
 10 & \ensuremath{3.8798\times 10^{-6}} & \ensuremath{3.59399\times 10^{-7}} & \ensuremath{3.16225\times 10^{-6}} \\
\hline
\end{tabular}
\end{table}


\begin{figure}[h]
  \centering
  \includegraphics[width=0.85\linewidth]{stencil_flow_leibniz_last_run.png}
  \caption{Leibniz-gated stencil selection under refinement. The gate filters local stencil operators by requiring a shrinking Leibniz defect under refinement, producing a concentrated set of derivative-like closures (see Table~\ref{tab:stencil-leibniz}).}
  \label{fig:stencil-leibniz}
\end{figure}

\subsection{Integration as a packaged closure of sums}
The same closure lens applies to integration: cumulative sums are staged at grid size $h$, audited by differentiation and protocol mismatch defects, and accepted as packaged closures when those defects shrink under refinement.
In the recorded run, the left-versus-trapezoid route mismatch at the finest grid is $\IntegrationRMAtSmallestH$, with fitted scaling exponent $\IntegrationRMExponent$ and fit quality $\IntegrationFitRtwo$.
The fundamental-theorem defect for the trapezoid cumulative operator decays with fitted exponent $\IntegrationFTExponent$.
(Here the reported norms are unweighted discrete $\ell^2$ norms over grid points, so pointwise $O(h)$ defects appear as $O(h^{1/2})$ in $\ell^2$.)

% Auto-generated by scripts/export_results_tex.py from notes/*_last_run.json
\begin{table}[h]
\centering
\caption{Integration closure diagnostics under refinement}
\label{tab:integration-closure}
\begin{tabular}{l|l|l|l|l}
\hline
 N & $h$ & RM & FT left & FT trap \\
\hline
 128 & \ensuremath{0.0078125} & \ensuremath{0.0141732} & \ensuremath{5.65355\times 10^{-14}} & \ensuremath{0.19633} \\
\hline
 256 & \ensuremath{0.00390625} & \ensuremath{0.00708551} & \ensuremath{1.62236\times 10^{-13}} & \ensuremath{0.138837} \\
\hline
 512 & \ensuremath{0.00195312} & \ensuremath{0.00354262} & \ensuremath{4.75038\times 10^{-13}} & \ensuremath{0.0981742} \\
\hline
 1024 & \ensuremath{0.976562\times 10^{-3}} & \ensuremath{0.00177129} & \ensuremath{1.2982\times 10^{-12}} & \ensuremath{0.0694199} \\
\hline
\end{tabular}
\end{table}


\begin{figure}[h]
  \centering
  \includegraphics[width=0.85\linewidth]{integration_closure_last_run.png}
  \caption{Integration closure diagnostics on staged grids. The route mismatch between left and trapezoid cumulative protocols and the differentiation-based defect both decay under refinement, consistent with stable packaged closure.}
  \label{fig:integration-closure}
\end{figure}
